% =====================================================================
% Section IV -- SpecAsync-UVM: design and the first-touch oracle
% SpecAsync-UVM manuscript, draft 1 (2026-10-10)
%
% SCOPE: what was built and how it can be configured. No results here
% (Sections VI-VII); numbers below are design constants or come from the
% verification gates that established correctness (E0.9, E3b).
%
% SOURCES:
%   hook, enqueue, queue cap, worker, trylock, THROTTLED:
%     uvm_gpu_replayable_faults.c (driver/src), E1B_HANDOFF_COST.md §1,
%     E3A_WIDTH_REVIEW.md §1-3, GATE_T4_REPORT.md §2 (cap 1024).
%   worker never maps: E0.9a-2 §3 (uvm_va_block.c:4968-4995, :4886-4887).
%   width W: E3A_WIDTH_REVIEW.md (power of two in [1,512], insmod rejects,
%     region bounds, liveness re-check, W=1 identical to pre-E3a).
%   first-touch table and cursor: GATE_E0_9_REPORT.md §3 (algorithm:
%     monotonic cursor, skip-correction, lookahead L, held, exhausted,
%     unknown page; sort at load, bsearch per fault; distinct-page assert).
%   fast lookup: GATE_E3B_REPORT.md (ft_fast verified at load; nranges
%     Stencil 2, GraphBFS 19; cas_giveup 0), E3B_ORACLE_COST.md.
%   identities I1, I2: GATE_E3B_REPORT.md Steps 2-3 (exact on every run).
%   E3a-2 = verified at W=1, fast=0: GATE_E3B_REPORT.md Step 2 (b) vs (a).
%   fault_already_resident: E0.9a-2 (counter on the demand path); CS2-N7.
%   decomposition overhead: CS2-16 (+0.331% 5070 Ti; +0.64% T4 UNCHECKED).
%   review protocol: E3A_WIDTH_REVIEW.md §3, §5-6; catalog #22.
%
% OPEN -- design checks (read-only, Claude Code):
%   [DC-1] same-region suppression: E3a used "last enqueued region"; the
%          handoff says E3a-2 has an 8-slot region history. Confirm which
%          the E3a-2 module (33FD42E6) implements and its exact rule.
%   [DC-2] ft_fast mechanism: confirm "rank computed from a small set of
%          contiguous address runs" and "cursor advanced by compare-and-swap
%          with bounded retries (cas_giveup)"; state what happens when the
%          load-time verification fails (fast disabled? load refused?).
%   [DC-3] table collection: confirm it runs with ASLR disabled (setarch -R)
%          and records once per coalesced fault after E0.5.
%   [DC-4] depth semantics: depth 0 = lookup only (no make_resident),
%          depth 1 = make_resident. Confirm.
%   [FIG]  Fig. 2: fault path with the hook, queue, worker and table.
% =====================================================================

\section{SpecAsync-UVM}
\label{sec:design}

SpecAsync-UVM adds an off-path staging mechanism to NVIDIA's open-source UVM
driver. Its purpose in this study is narrow: to let us measure what moving
residency work off the fault-servicing thread can and cannot achieve. We
therefore describe it as an instrument, with the parameters we vary and the
checks that establish it does what it claims.

\subsection{Structure}
\label{sec:design-structure}

The mechanism has three parts (Fig.~2).

\paragraph{A prediction hook on the servicing thread}
At the start of servicing each batch, before the address-space lock is
taken, the driver walks the batch's coalesced faults and asks a
\emph{policy} for one predicted page per fault. A policy may decline. When it
returns a page, the hook allocates a small work item, takes a reference on the
GPU, and queues the item on a kernel workqueue. The queue is capped at 1{,}024
outstanding items; a prediction that arrives when the queue is full is
dropped and counted. This hook is the only part of the mechanism that runs on
the fault-servicing thread, and it runs on every fault, whether or not a
prediction results. Section~VI measures its cost.

\paragraph{A worker off the critical path}
A kernel worker thread takes items from the queue. For each, it takes the
address space's lock in read mode, finds the VA block containing the
predicted page, and \emph{tries} to take that block's lock. It never waits:
if the block is locked, typically because the servicing thread is working on
it, the item is abandoned and counted as throttled. With the lock held, it
calls the driver's \texttt{make\_resident} routine with the predicted region,
the GPU as destination, and a prefetch cause. The worker therefore performs
exactly the operation that Section~\ref{sec:bg-resident-mapped} showed cannot
install a GPU mapping. This is a property of the driver's interfaces, not a
choice we made to weaken the mechanism: the mapping step is part of fault
service, and invoking it from outside a fault would require restructuring the
driver's service path, which we leave to future work.

\paragraph{Policies}
A policy maps the fault being serviced to a page worth staging. Earlier
phases of this project evaluated adjacent-page and stride heuristics and a
trace-replay oracle; the last of these was misaligned
(Section~VIII) and none of their results is used in this paper. Every
speculation result reported here uses one policy, the first-touch oracle of
Section~\ref{sec:design-oracle}, against a control with speculation disabled.

A \emph{depth} parameter selects how much of the worker's work is done:
depth~0 performs the lookup without migrating, and depth~1 migrates. All
reported configurations use depth~1.

\subsection{Staging Width}
\label{sec:design-width}

At one page per item the worker cannot keep pace with demand: with a fast
predictor the queue overflows, and each item pays the full cost of a lock
acquisition and a migration call for 4\,KB of data. A \emph{width} parameter
$W$ addresses this. For a predicted page at index $i$ in its VA block, the
worker stages the aligned region
$[\,\lfloor i/W \rfloor W,\; \min(\lfloor i/W \rfloor W + W,\; n)\,)$, where $n$
is the number of pages in the block; at $W = 512$ the region is the whole
2\,MB block. On the prediction side, a prediction whose region was already
enqueued recently is not enqueued again, so that one item covers a region
rather than one item per page of it. % [DC-1] exact suppression rule.

$W$ is a load-time module parameter, accepted only as a power of two from 1
to 512; any other value makes the module load fail. Every region bound is
checked explicitly before the worker builds its page mask, because the
driver's own helpers check some of them only with assertions that are
non-fatal in release builds, and one of them writes the mask without any
check. The worker also re-checks, under the block lock, that the block has
not been destroyed. At $W = 1$ every new code path is bypassed and the
mechanism behaves exactly as it did before width was added; we verified this
counter for counter against the previously verified module.

\subsection{The First-Touch Oracle}
\label{sec:design-oracle}

To separate what off-path staging can achieve from how well a heuristic
predicts, we give the worker knowledge of the future taken from the workload
itself. In a collection run with address-space randomisation disabled, the
driver records each coalesced fault's page in service order. A user-space
tool reduces this record to its \emph{first-touch order}: the list of
distinct pages in the order each was first faulted, ranked
$0, 1, \ldots, N-1$. The table is loaded into the module before a measured
run. % [DC-3]

During the measured run the policy keeps a cursor $c$, the rank of the next
page it may predict. On a fault to a page of rank $r$:
\begin{itemize}
  \item if $c \le r$, demand has caught up with the cursor, which skips to
        $r + 1$;
  \item if $c \le r + L$, the policy predicts the page of rank $c$ and
        advances the cursor; otherwise it holds back, so that it never runs
        more than $L$ ranks ahead of demand;
  \item a fault to a page absent from the table, or a cursor past the end of
        the table, produces no prediction.
\end{itemize}
Each table entry is predicted at most once per run, and the lookahead $L$
bounds how far ahead staging can run.

This oracle is as good as the workload's run-to-run stability allows, and
no better. On our stencil workload, the first-touch order of 1{,}125{,}000
pages differs between independent runs by at most 107 ranks, so the table is
a near-perfect forecast. On graph traversal the order is much looser, and
because the cursor skips forward whenever a fault arrives ahead of it, one
out-of-order fault can jump the cursor past pages that have not yet been
touched. The graph oracle therefore stages only part of the working set, and
we report its results as those of a weakened oracle. An index-based replay
of the raw fault stream, which we used in earlier phases, is not robust in
the same way: the number of repeated faults per page varies from run to run,
so positions drift even when the order of first touches does not.

\paragraph{The lookup must be cheap}
The original implementation found a page's rank by binary search over a
sorted copy of the table and updated the cursor under a global spinlock with
interrupts disabled, once for every coalesced fault. That cost falls on the
servicing thread, so an oracle meant to bound speculation's benefit was
itself reducing it. A second implementation computes the rank in constant
time and advances the cursor without a lock. At load time the module checks
the fast lookup against the original for every entry in the table, and
records whether the check passed; every measured run with the fast lookup
requires it to have passed. % [DC-2] mechanism and failure behaviour.

\subsection{Telemetry}
\label{sec:design-telemetry}

The module exports counters through debugfs and records per-batch timing in
fixed-size rings that drop rather than overwrite when full, and count their
drops. Three kinds of record matter for the results.

\paragraph{Accounting identities}
Every coalesced fault is either predicted, held back, unknown, past the end
of the table, or abandoned by the lock-free update; and every prediction is
either suppressed as a repeat region, enqueued, or dropped at a full queue.
Both identities are checked on every run, and held exactly on every run we
report.

\paragraph{Phase timing}
Timestamps at the boundaries of Table~\ref{tab:phases} give the
decomposition of each batch, together with the time the servicing thread
spends in the prediction hook before the address-space lock. On the
RTX~5070~Ti the decomposition instrumentation adds 0.33\% to kernel-loop time,
which is not distinguishable from zero.

\paragraph{Whether staging arrived first}
A counter on the demand path records each fault whose page was already
resident on the GPU when the servicing thread reached it, that is, faults
for which the worker finished first. It measures that the worker won the
race to make the page resident, not that the fault was avoided, since by
Section~\ref{sec:bg-resident-mapped} it was not. An older counter of
``hits'', faults on pages predicted within the previous 10\,ms, is not a
measure of success and is not used.

\subsection{Developing Kernel Code Safely}
\label{sec:design-safety}

The mechanism runs inside a production kernel module, where a mistake can
crash the machine. One did: an instrumentation counter added during this
study called a driver helper that returns a null pointer when the GPU has no
state in a block, a case the helper guards only with a non-fatal assertion,
and the kernel oopsed. Since then every change to kernel code follows the
same protocol. The change is compiled but not loaded. A review document lists
every driver function it calls, with the locks each requires, whether it can
fail or return null, and every precondition that is checked only by
assertion, and states how the new code guards each one. Invalid parameter
values are rejected at load. The first load is attended, starts with the
smallest probe, and checks the new module against the previously verified
one at settings where they should be identical. The verified module is never
overwritten, and every measured run checks the loaded module's source hash
and parameters before it starts.

\subsection{Configurations}
\label{sec:design-configs}

Table~\ref{tab:configs} lists the configurations used in Sections~VI
and~VII. C0 is the comparison of record throughout: the stock prefetcher at
its shipped threshold, with no speculation. C1 disables the prefetcher and is
used only to isolate mechanisms, never as the baseline against which
speculation is judged.

\begin{table}[t]
\centering
\footnotesize
\setlength{\tabcolsep}{3.5pt}
\begin{tabular}{@{}l l l l l@{}}
\toprule
\textbf{Config} & \textbf{Speculation} & \textbf{Prefetcher} & $W$ & \textbf{Lookup} \\
\midrule
C0        & off                      & on (threshold $t$)  & --- & --- \\
C1        & off                      & off                 & --- & --- \\
C6-W$k$   & first-touch oracle       & off                 & $k$ & fast \\
C7-W$k$   & first-touch oracle       & on (threshold $t$)  & $k$ & fast \\
\bottomrule
\end{tabular}
\caption{Configurations. Unless stated, $t = 51$ (shipped), lookahead
$L = 4096$, depth~1. ``Slow'' variants use the original binary-search lookup.
Stock-$t$ denotes the unmodified driver at threshold $t$.}
\label{tab:configs}
\end{table}

% [FIG] Fig. 2: Section III's fault-path figure extended with: the prediction
% hook before the address-space lock (on the servicing thread), the bounded
% queue, the worker (VA-space read lock -> block trylock -> make_resident,
% no map), the first-touch table with cursor c and lookahead L.
